% Source: https://artofproblemsolving.com/community/c3h1735512_2019_mathcounts_handbook_warmup_6
% Stable identity: aops-topic-1735512-q82
Every line through $A$ meets circle $A$, since $A$ is its centre. So we need the probability that the line meets circle $B$.

Let the square have side $s$. Then circle $B$ has radius $BE=s$, and $AB=s\sqrt2$ is a diagonal of the square. Since $\angle AEB=90^\circ$, the line $AE$ is perpendicular to the radius $BE$, so it is tangent to circle $B$; likewise $AF$. These tangents make angles of $45^\circ$ with $AB$ (the diagonal of a square bisects its angle):

\begin{asy}
usepackage("amsmath");
size(6cm);
defaultpen(fontsize(10pt));
real s=2; pair A=(0,0),B=(s*sqrt(2),0),Ep=(s/sqrt(2),s/sqrt(2)),F=(s/sqrt(2),-s/sqrt(2));
fill(A--(4.4,4.4)--(4.4,-4.4)--cycle,rgb(1,0.92,0.85));fill(A--(-1.2,1.2)--(-1.2,-1.2)--cycle,rgb(1,0.92,0.85));
draw(circle(A,s),linewidth(0.8));draw(circle(B,s),linewidth(0.8));
draw(A--Ep--B--F--cycle,blue);
draw((-1.2,-1.2)--(4.4,4.4),red+linewidth(0.8));draw((-1.2,1.2)--(4.4,-4.4),red+linewidth(0.8));
draw(A--B,dashed);
draw(Ep+0.3*dir(-45)--Ep+0.3*dir(-45)+0.3*dir(-135)--Ep+0.3*dir(-135));
dot(A);dot(B);
label("$A$",A+(-0.5,0));label("$B$",B,E);label("$E$",Ep,2*N);label("$F$",F,2*S);
label("\mbox{}",(min(currentpicture,user=true)+max(currentpicture,user=true))/2,black); // keeps the text after the figure black when LaTeX typesets the labels inline
\end{asy}

A line through $A$ meets circle $B$ exactly when its angle with line $AB$ is at most $45^\circ$ (the shaded double wedge bounded by the two tangent lines). A random line through $A$ has its direction uniformly distributed over $180^\circ$, and $90^\circ$ of these directions work, so the probability is $\dfrac{90}{180}=\boxed{\dfrac{1}{2}}$.
